\documentclass[11pt,reqno]{amsart}
\usepackage[T1]{fontenc}
\usepackage{lmodern,amssymb,microtype}
\usepackage[colorlinks=true,linkcolor=black,citecolor=blue,urlcolor=black]{hyperref}
\hypersetup{pdftitle={A counterexample to the Erdos--Mahler conjecture},
  pdfsubject={Smooth convergents and prescribed irrationality exponents},
  pdfauthor={Diego Marques}}

\usepackage{enumitem}

\newtheorem*{theorem}{Theorem}
\newtheorem{lemma}{Lemma}[section]
\newtheorem{proposition}[lemma]{Proposition}
\theoremstyle{definition}
\newtheorem*{conjecture}{Conjecture}
\newcommand{\N}{\mathbb N}
\newcommand{\Q}{\mathbb Q}
\newcommand{\Z}{\mathbb Z}
\newcommand{\DD}{\mathcal D}
\newcommand{\PP}{\mathrm P}
\setlength{\emergencystretch}{1em}

\title{A counterexample to the Erd\H{o}s--Mahler conjecture}
\author[D. Marques]{Diego Marques}
\address{Departamento de Matem\'atica, Universidade de Bras\'ilia,
Bras\'ilia, DF, Brazil}
\email{diego@mat.unb.br}
\thanks{The author was supported by the Brazilian National Council for
Scientific and Technological Development (CNPq), Grant No.~304467/2023-5.}
\subjclass[2020]{Primary 11J70; Secondary 11J82, 11J86, 11J87}
\keywords{Erd\H{o}s--Mahler conjecture, continued fractions, smooth
convergents, irrationality exponent, Subspace Theorem, linear forms in
logarithms}

\begin{document}

\begin{abstract}
We disprove a 1939 conjecture of Erd\H{o}s and Mahler by constructing,
for every real $\tau\ge4$, a Cantor set of numbers of irrationality
exponent $\tau$, each having infinitely many regular continued-fraction
convergents of the form $2^a/3^b$.
\end{abstract}
\maketitle

\section{Introduction}\label{sec:intro}

For an irrational real number $\xi$, its \emph{irrationality exponent}
$\mu(\xi)$ is the supremum of the real numbers $\nu>0$ for which
$0<|\xi-u/v|<v^{-\nu}$ holds for infinitely many reduced fractions
$u/v$ with $v\ge1$. Thus $\xi$ is a Liouville number precisely when
$\mu(\xi)=\infty$. Write $p_n/q_n$ for the regular continued-fraction
convergents of $\xi$, and $\PP(m)$ for the largest prime factor of
$|m|$ when $m$ is a nonzero integer, with $\PP(\pm1)=1$.

In 1939, Erd\H{o}s and Mahler~\cite{EM} investigated
the prime factors of continued-fraction numerators
and denominators. They proved that if
$\PP(q_{n-1}q_nq_{n+1})$ is bounded for infinitely
many $n$, then $\xi$ is a Liouville number.
They conjectured that the same conclusion should
hold if $\PP(p_nq_n)$ is bounded for infinitely
many $n$. The question was subsequently studied by
Fraenkel~\cite{Fraenkel} and recalled by Erd\H{o}s
in his reminiscences of Mahler~\cite{ErdosReminiscences}. 
In 2004, Bugeaud~\cite[Problem~43, p.~215]{Bugeaud}
included the conjecture among the open problems
on continued fractions.

\begin{conjecture}[Erd\H{o}s--Mahler]
If $\PP(p_nq_n)$ is bounded for infinitely many $n$, then $\xi$ is
a Liouville number.
\end{conjecture}

We disprove the conjecture even when the selected numerators are
powers of $2$ and the corresponding denominators are powers of $3$.

\begin{theorem}
For every real $\tau\ge4$, there exists a Cantor set
$K_\tau\subset(1,2)$ such that every $\xi\in K_\tau$
has irrationality exponent $\tau$ and infinitely many regular
continued-fraction convergents of the form $2^a/3^b$, with $a,b$
positive integers.
\end{theorem}

By Ridout's theorem~\cite{Ridout}, bounded prime factors of the
denominators of infinitely many convergents already imply
transcendence. Fraenkel~\cite{Fraenkel62,Fraenkel} studied rational
approximation with prescribed prime factors and constructed Liouville
numbers with infinitely many convergents subject to such restrictions.
Fraenkel and Borosh~\cite{FraenkelBorosh} subsequently obtained
non-Liouville examples by allowing additional cofactors in the
numerators and denominators. When both cofactors are bounded, their
Hausdorff-dimension formula gives zero. Lelis and Marques~\cite[Theorem~1.1]{LelisMarques}
proved the conjectured conclusion under an additional
spacing condition on the indices of the selected convergents.

More recently, Bugeaud and Nguyen~\cite[Question~4.1, p.~32]{BN}
revisited the conjecture and constructed further examples with
restricted prime factors, leaving open whether the numbers they
constructed are Liouville. This illustrates the main difficulty:
prescribing infinitely many convergents does not exclude
exceptionally good rational approximations at other denominators.


The key ingredient is a counting estimate for the
prescribed fractions that admit exceptionally close
rational approximations. The estimate is independent
of the intervals used in the construction and, together
with the distribution of an irrational rotation, ensures
that sufficiently many choices remain at each stage.
This allows us to construct Cantor sets with the
prescribed irrationality exponent.

The paper is organized as follows. Section~2 states the
counting proposition and deduces the main theorem from it.
Section~3 proves the proposition, using rotation estimates,
the $p$-adic Subspace Theorem, and multiplicative relations
among the prescribed fractions.


\section{The counting proposition and the main theorem}
\label{sec:main}

We first state the counting estimate underlying the
construction and deduce the main theorem from it.
The proof of the estimate is given in the next section.

\smallskip
\noindent\emph{Notation.}
All logarithms are natural, and $\N=\{1,2,\ldots\}$. Rational
fractions have positive denominators. We write $|I|$ for the
length of an interval or arc and $|A|$ for the cardinality of a
finite set. The symbols $O$ and $\ll$ have their usual meaning;
implied constants are absolute unless a dependence is indicated.

Set
\[
 \alpha=\frac{\log3}{\log2},\qquad
 a(b)=\lceil b\alpha\rceil,\qquad
 R(b)=\frac{2^{a(b)}}{3^b}\qquad(b\in\N).
\]
Unique factorization shows that $\alpha$ is irrational, so
$0<a(b)-b\alpha<1$. Consequently, $R(b)$ is the unique reduced
fraction in $(1,2)$ with denominator $3^b$ and numerator a power
of $2$. Let $\DD$ be the set of denominators of the regular
continued-fraction convergents of $\alpha$. For $x\in\mathbb R$
write $\{x\}=x-\lfloor x\rfloor$, and define
\[
 \Phi(y)=1-\frac{\log y}{\log2}\qquad(1\le y\le2).
\]
Then $\Phi(R(b))=\{b\alpha\}$ and
$|\Phi'(y)|=(y\log2)^{-1}$.

For $4<\nu\le5$, a reduced fraction $u/v\in[1,2]$ is
\emph{exceptional for $b$ at exponent $\nu$} if
$u/v\ne R(b)$ and
\begin{equation}\label{eq:exception}
 \left|R(b)-\frac uv\right|<\frac1{16v^\nu}.
\end{equation}
For $N\in\N$ and $D>0$, let $E_N(D,\nu)$ be the number of
integers $b\in[N,2N)$ admitting an exceptional fraction
with $v>D$.

\begin{proposition}[Exceptional approximations]
\label{prop:count}
There exists an absolute constant $C>0$ such that,
for every $4<\nu\le5$ and $D\ge2$,
\[
 \limsup_{\substack{N\to\infty\\N\in\DD}}
 \frac{E_N(D,\nu)}{N}\le CD^{2-\nu}.
\]
\end{proposition}

The deduction of the theorem also uses an elementary
estimate for the distribution of the fractions $R(b)$,
obtained from the convergents of $\alpha$.

\begin{lemma}[Rotation estimate]\label{lem:rotation}
For every arc $I\subset\mathbb R/\mathbb Z$ and every
$N\in\DD$,
\begin{equation}\label{eq:grid}
 \#\{b\in\Z:N\le b<2N,\ \{b\alpha\}\in I\}
 =N|I|+O(1),
\end{equation}
with an absolute implied constant.
\end{lemma}

\begin{proof}
Choose a convergent $p/N$ of $\alpha$. Since
$\gcd(p,N)=1$, the points
\[
 N\alpha+\frac{ip}{N}\pmod1,\qquad 0\le i<N,
\]
form a translate of the uniform grid of spacing $1/N$.
Thus $N|I|+O(1)$ of these points lie in $I$.
By the convergent estimate
$|\alpha-p/N|<N^{-2}$
\cite[Corollary~1.4, p.~10]{Bugeaud},
\[
 \left|(N+i)\alpha-
 \left(N\alpha+\frac{ip}{N}\right)\right|
 =i\left|\alpha-\frac pN\right|<\frac1N.
\]
The two counts can therefore differ only at grid points
within distance $1/N$ of an endpoint of $I$. There are
$O(1)$ such points, uniformly in $I$ and $N$, proving
\eqref{eq:grid}.
\end{proof}

We turn to the construction. The rotation estimate
and Proposition~\ref{prop:count} ensure that two
admissible choices remain at each stage.

Fix $\tau\ge4$. For $t\in\N$, put $Q=3^t$ and define
\[
 \begin{aligned}
 J_t&=(R(t)+Q^{-\tau},\,R(t)+2Q^{-\tau}),\\
 K(t)&=[R(t)+Q^{-\tau},\,R(t)+3Q^{-\tau}].
 \end{aligned}
\]
Since $R(t)\le2-Q^{-1}$, both intervals lie in $(1,2)$
for sufficiently large $t$. Set
\[
 D_t=\frac{Q^{\tau-1}}6,\qquad
 \nu_t=4+\frac1t.
\]
An index $b$ is \emph{admissible at $t$} if
$R(b)\in J_t$ and no fraction exceptional for $b$
at exponent $\nu_t$ has denominator greater than $D_t$.

\begin{lemma}[Admissible choices]\label{lem:choice}
For each $\tau\ge4$, there exists $t_0$ such that,
for every integer $t\ge t_0$, the interval $[N,2N)$
contains at least two distinct indices admissible at $t$
whenever $N\in\DD$ is sufficiently large, depending on $t$.
\end{lemma}

\begin{proof}
Fix $t$ sufficiently large that $J_t\subset(1,2)$
and $D_t\ge2$, and put $\rho_t=|\Phi(J_t)|$.
Since
\[
 \frac1{2\log2}\le|\Phi'(y)|
 \le\frac1{\log2}\qquad(1\le y\le2),
\]
the mean value theorem gives
\[
 \frac{Q^{-\tau}}{2\log2}
 \le\rho_t\le\frac{Q^{-\tau}}{\log2}.
\]

By the rotation estimate \eqref{eq:grid}, there are
$N\rho_t+O(1)$ indices $b\in[N,2N)$ for which
$R(b)\in J_t$. Of these, at most $E_N(D_t,\nu_t)$
are inadmissible. If $A_t(N)$ denotes the number
of admissible indices, Proposition~\ref{prop:count}
therefore yields
\[
 \liminf_{\substack{N\to\infty\\N\in\DD}}
 \frac{A_t(N)}N
 \ge\rho_t-CD_t^{2-\nu_t}
 \ge\rho_t-36C Q^{-2(\tau-1)},
\]
where $C$ is the absolute constant in that proposition
and we used $\nu_t>4$.

The error relative to $\rho_t$ satisfies
\[
 \frac{36C Q^{-2(\tau-1)}}{\rho_t}
 \le72C(\log2)Q^{2-\tau}\longrightarrow0
 \qquad(t\to\infty),
\]
since $\tau\ge4$. Hence, for all sufficiently large
$t$, the preceding lower limit is at least
$\rho_t/2>0$. With such a $t$ fixed, $A_t(N)$
tends to infinity as $N\to\infty$ through $\DD$.
In particular, there are at least two admissible
indices for every sufficiently large $N\in\DD$.
\end{proof}

\begin{proof}[Proof of the theorem]
Choose $t_0$ sufficiently large for Lemma~\ref{lem:choice}.
For each $t\ge t_0$, take $N\in\DD$ with
$N>Q=3^t$, large enough to provide two admissible
indices $b<c$ in $[N,2N)$. We first show that
$K(b)$ and $K(c)$ are disjoint and lie in the
interior of $K(t)$.

For $d=b,c$, admissibility gives $R(d)\in J_t$.
Moreover, $d>t$, so $3^d\ge2Q$. The left endpoint
of $K(d)$ is greater than $R(t)+Q^{-\tau}$,
whereas its right endpoint satisfies
\[
 \begin{aligned}
 R(d)+3\cdot3^{-\tau d}
 &<R(t)+(2+3\cdot2^{-\tau})Q^{-\tau}\\
 &<R(t)+3Q^{-\tau},
 \end{aligned}
\]
since $\tau\ge4$. Thus
\[
 K(b),K(c)\subset\operatorname{int}K(t).
\]

To prove disjointness, observe that $c<2b$.
The distinct reduced fractions $R(b)$ and $R(c)$
can both be expressed with denominator $3^c$.
Consequently,
\[
 |R(b)-R(c)|
 \ge3^{-c}>3^{-2b}
 >3\cdot3^{-\tau b}.
\]
Each of $K(b)$ and $K(c)$ lies to the right
of its defining fraction, at distance at most
$3\cdot3^{-\tau b}$. Hence the two intervals
are disjoint.

Fix such a pair of successors for every $t\ge t_0$
and iterate the construction, labeling the two
choices by $0$ and $1$. Each infinite binary
sequence determines a nested chain
\[
 K(t_0)\supset K(t_1)\supset K(t_2)\supset\cdots.
\]
Along this chain, put
$r_j=R(t_j)$ and $Q_j=3^{t_j}$. Since
$Q_{j+1}\ge2Q_j$, the intervals satisfy
\[
 |K(t_j)|=2Q_j^{-\tau}
 \le2Q_0^{-\tau}2^{-\tau j}.
\]
Their intersection therefore consists of a unique
point $\xi$. The uniform decay of the interval
lengths makes the resulting map from
$\{0,1\}^{\N}$ continuous, while disjointness
of the successors makes it injective.
By compactness, this map is a homeomorphism
onto its image, a Cantor set
$K_\tau\subset(1,2)$.

Fix $\xi\in K_\tau$ and consider its associated
chain. Strict nesting gives
\begin{equation}\label{eq:displacement}
 Q_j^{-\tau}<\xi-r_j<3Q_j^{-\tau}
 \qquad(j\ge0).
\end{equation}
We first show that $\xi$ is irrational.
If $\xi=u/v$ were rational, then
$\xi\ne r_j$ and rational separation would give
\[
 |\xi-r_j|\ge\frac1{vQ_j}.
\]
Since $Q_j\to\infty$ and $\tau\ge4$, this
contradicts \eqref{eq:displacement} for large $j$.

Moreover, for every $j$,
\[
 0<\xi-r_j<3Q_j^{-\tau}
 <\frac1{2Q_j^2}.
\]
As each $r_j$ is reduced, Legendre's criterion
\cite[Theorem~1.8, p.~10]{Bugeaud} shows that
$r_j$ is a regular continued-fraction convergent
of $\xi$. These convergents have numerators
that are powers of $2$ and strictly increasing
denominators that are powers of $3$. For every
$0<s<\tau$, \eqref{eq:displacement} gives
\[
 0<\xi-r_j<Q_j^{-s}
\]
for all sufficiently large $j$. Hence
$\mu(\xi)\ge\tau$.

We now prove the reverse inequality. Let
$u/v\in[1,2]$ be reduced with $v\ge Q_0$,
and choose the unique $j$ such that
\[
 Q:=Q_j\le v<Q_{j+1}=:S.
\]
If $u/v=r_j$, then $v=Q$ and
\eqref{eq:displacement} gives
$|\xi-u/v|>v^{-\tau}$. We may therefore
assume $u/v\ne r_j$ and distinguish two
ranges of denominators.

\smallskip
\noindent\emph{(i) $v\le D_{t_j}$.}
Since $r_j$ is reduced with denominator $Q$,
rational separation and \eqref{eq:displacement}
give
\[
 \begin{aligned}
 \left|\xi-\frac uv\right|
 &\ge\frac1{vQ}-\frac3{Q^\tau}\\
 &\ge\frac1{2vQ}
 \ge\frac1{2v^2},
 \end{aligned}
\]
where the second inequality follows from
$6v\le Q^{\tau-1}$ and the last from $Q\le v$.

\smallskip
\noindent\emph{(ii) $v>D_{t_j}$.}
Since $v<S$ and $r_{j+1}$ is reduced with
denominator $S$, we have $u/v\ne r_{j+1}$ and
\[
 \left|r_{j+1}-\frac uv\right|\ge\frac1{Sv}.
\]
On the other hand, $S\ge3$ and $v<S$ imply
$6v<S^3$. As $\tau\ge4$,
\eqref{eq:displacement} gives
\[
 \begin{aligned}
 |\xi-r_{j+1}|
 &<\frac3{S^\tau}
 \le\frac3{S^4}
 <\frac1{2Sv}\\
 &\le\frac12\left|r_{j+1}-\frac uv\right|.
 \end{aligned}
\]
The index $t_{j+1}$ is admissible at $t_j$.
Since $v>D_{t_j}$, the definition of
admissibility yields
\[
 \left|r_{j+1}-\frac uv\right|
 \ge\frac1{16v^{\nu_{t_j}}}.
\]
Consequently,
\[
 \left|\xi-\frac uv\right|
 \ge\frac12\left|r_{j+1}-\frac uv\right|
 \ge\frac1{32v^{\nu_{t_j}}}.
\]

Combining the two ranges, we have
\[
 \left|\xi-\frac uv\right|
 \ge\frac1{32v^{\nu_{t_j}}}
 \qquad
 (Q_j\le v<Q_{j+1},\ u/v\ne r_j).
\]
Since $t_j\to\infty$, we have
$\nu_{t_j}\to4$. Given $\varepsilon>0$,
it follows that $\nu_{t_j}\le4+\varepsilon/2$
for all sufficiently large $j$. Moreover,
$v^{\varepsilon/2}\ge32$ for all sufficiently
large $v$. Thus
\[
 \left|\xi-\frac uv\right|
 \ge v^{-4-\varepsilon}
 \ge v^{-\tau-\varepsilon}
\]
for every sufficiently large denominator $v$
and every reduced $u/v\in[1,2]$ distinct
from the prescribed convergent in its range.

For the prescribed convergents,
\eqref{eq:displacement} gives
$|\xi-r_j|>Q_j^{-\tau}$.
Finally, rationals outside $[1,2]$ are
bounded away from $\xi$, while only finitely
many reduced fractions in $[1,2]$ have
bounded denominator. It follows that
$\mu(\xi)\le\tau$. Together with the lower
bound, this proves
\[
 \mu(\xi)=\tau
 \qquad\text{for every }\xi\in K_\tau.
\]
\end{proof}

\section{Proof of the counting proposition}
\label{sec:count}

We begin with a direct estimate for small exceptional denominators
and a height bound for the remaining ones. We then count the
exceptional fractions in fixed denominator ranges by exploiting
the multiplicative structure of $R(b)$, and sum the resulting estimates.

\subsection{Small denominators and a height bound}

The rotation estimate gives a direct count for exceptional fractions
with denominator at most $N^{1/3}$. The $p$-adic Subspace Theorem
will provide the height bound needed for larger denominators.

\begin{lemma}[Small-denominator approximations]
\label{lem:small}
There is an absolute constant $C>0$ such that, for every
$4<\nu\le5$, $D\ge2$, and $N\in\DD$,
\begin{equation}\label{eq:small}
 E_N(D,\nu)\le E_N(N^{1/3},\nu)
 +C\bigl(ND^{2-\nu}+N^{2/3}\bigr).
\end{equation}
\end{lemma}

\begin{proof}
If $D\ge N^{1/3}$, the assertion follows from the
monotonicity of $E_N(D,\nu)$ in $D$. Suppose therefore
that $D<N^{1/3}$.

For $V\ge2$, there are $O(V^2)$ reduced fractions
$u/v\in[1,2]$ with $V<v\le2V$, since each denominator
$v$ admits at most $v+1$ numerators. As
$|\Phi'(y)|\le(\log2)^{-1}$ on $[1,2]$, the image
under $\Phi$ of the $(16v^\nu)^{-1}$-neighborhood of
any such fraction, intersected with $[1,2]$, is an
interval of length $O(V^{-\nu})$. By
$\Phi(R(b))=\{b\alpha\}$ and \eqref{eq:grid}, this
fraction can be exceptional for at most
$O(NV^{-\nu}+1)$ indices $b\in[N,2N)$.

Thus denominators in $(V,2V]$ contribute at most
$O(NV^{2-\nu}+V^2)$ indices. Cover $(D,N^{1/3}]$
by the intervals $(V,2V]$, where $V=2^iD$ and
$i\ge0$, truncating the last one if necessary.
Their total contribution is at most a constant times
\[
 \sum_{\substack{i\ge0\\2^iD<N^{1/3}}}
 \bigl(N(2^iD)^{2-\nu}+(2^iD)^2\bigr)
 \ll ND^{2-\nu}+N^{2/3}.
\]
Indeed, the first sum is geometric with ratio
$2^{2-\nu}\le1/4$, whereas the second is dominated
by its largest term. The implied constants are
absolute for $4<\nu\le5$. Adding the indices
admitting exceptional denominators greater than
$N^{1/3}$ gives \eqref{eq:small}.
\end{proof}

By Lemma~\ref{lem:small}, the proof of Proposition~\ref{prop:count}
reduces to showing that, for each fixed $4<\nu\le5$,
\[
 E_N(N^{1/3},\nu)=o(N)
 \qquad(N\to\infty,\ N\in\DD).
\]

For distinct reduced fractions $R(b)$ and $u/v$, rational
separation gives
\[
 \left|R(b)-\frac uv\right|\ge\frac1{3^bv}.
\]
Together with \eqref{eq:exception}, this yields
$16v^{\nu-1}<3^b$, and hence only $\log v=O(N)$ when
$N\le b<2N$. The summation over larger denominator ranges
requires a bound $\log v=o(N)$, uniform over the exceptional
fractions in each block.

Fix $4<\nu\le5$, and let $W_N$ be the largest value of
$\log v$ among fractions exceptional for some
$b\in[N,2N)$, with $W_N=0$ if there are none. The
elementary bound just obtained ensures that $W_N$ is finite.

\begin{lemma}[Height bound]\label{lem:height}
For every $\varepsilon>0$ there exists $b_0(\varepsilon)$
such that, for every $b\ge b_0(\varepsilon)$ and every
reduced fraction $u/v\in[1,2]$ distinct from $R(b)$,
\[
 \left|R(b)-\frac uv\right|
 \ge3^{-\varepsilon b}v^{-2}.
\]
Consequently, for each fixed $4<\nu\le5$,
\begin{equation}\label{eq:height}
 W_N=o(N)\qquad(N\to\infty,\ N\in\DD).
\end{equation}
\end{lemma}

\begin{proof}
Put $P=2^{a(b)}$ and $S=3^b$, so that $S<P<2S$.
For a reduced fraction $u/v\in[1,2]$ distinct from
$P/S$, set
\[
 d=\gcd(Pv,Su),\qquad
 U=\frac{Pv}{d},\qquad W=\frac{Su}{d},\qquad
 e=|Pv-Su|.
\]
Then $e\ge1$ and $(U,W)=1$. Every prime divisor
of $d$ is $2$ or $3$; otherwise it would divide
both $u$ and $v$. For a positive integer $n$,
write $[n]_6=2^{v_2(n)}3^{v_3(n)}$, with $v_p$
the $p$-adic valuation. Normalize the $p$-adic
absolute values by $|p|_p=p^{-1}$. Since $d$
has no prime divisors outside $\{2,3\}$,
\[
 |UW|_2|UW|_3=[UW]_6^{-1}
 =\frac{d^2}{PS[uv]_6}.
\]
Take the linear forms $X-Y,Y$ at the real place
and $X,Y$ at $2$ and $3$. Their product at $(U,W)$
satisfies
\[
 \Pi:=|(U-W)W|\,|UW|_2|UW|_3
 =\frac{eu}{P[uv]_6}
 \le2v^2\left|\frac PS-\frac uv\right|,
\]
where we used $e=Sv|P/S-u/v|$, $u\le2v$, and $S<P$.

Suppose that the first assertion fails for arbitrarily
large $b$. For the corresponding fractions, rational
separation gives
\[
 \frac1{Sv}\le\left|\frac PS-\frac uv\right|
 <S^{-\varepsilon}v^{-2},
\]
so $v<S$. Hence, with $H=\max(U,W)$, we have
$H\le2Sv<2S^2$, and therefore
\[
 \Pi<2S^{-\varepsilon}
 \le2^{1+\varepsilon/2}H^{-\varepsilon/2}.
\]
Let $\delta=\min(\varepsilon/4,1/2)>0$. For
sufficiently large $H$, the last bound is smaller
than $(U^2+W^2)^{-\delta/2}$. The pairs $(U,W)$
are primitive, and the two forms at each place
are linearly independent with determinant one.
The $p$-adic Subspace Theorem
\cite[p.~246, (1.3)--(1.5)]{Schlickewei} therefore
places all such pairs of sufficiently large height
on finitely many rational lines through the origin.
Each line contains at most one positive primitive
pair, and there are only finitely many pairs of
bounded height. Thus only finitely many pairs
$(U,W)$ can arise.

On the other hand,
\[
 0<\left|\frac UW-1\right|
 =\frac{|P/S-u/v|}{u/v}
 <S^{-\varepsilon}v^{-2}\longrightarrow0
 \qquad(b\to\infty),
\]
which is impossible for a finite set of positive
primitive pairs. This proves the first assertion.

For $N$ sufficiently large in terms of $\varepsilon$,
the first assertion and \eqref{eq:exception} give
$16v^{\nu-2}<3^{\varepsilon b}$ for every fraction
exceptional at exponent $\nu$ with $N\le b<2N$.
Since $b<2N$,
\[
 \limsup_{\substack{N\to\infty\\N\in\DD}}
 \frac{W_N}{N}
 \le\frac{2\varepsilon\log3}{\nu-2}.
\]
Letting $\varepsilon\to0$ proves \eqref{eq:height}.
\end{proof}

\subsection{Products in a fixed denominator range}

The key observation is that the products $R(b)R(c)$,
with $b,c\in[N,2N)$, take only $O(N)$ distinct values.
We show that this property persists for sufficiently
close rational approximants whose denominators lie
in a suitable range.

Fix $4<\nu\le5$, $1<\theta<\nu/4$, and $N\in\DD$.
For $w>0$, put
\[
 V=e^{w/\theta},\qquad H=e^w=V^\theta.
\]
Let $\mathcal A(w)$ be the set of indices $b\in[N,2N)$
admitting an exceptional fraction at exponent $\nu$
with denominator
\[
 \max\{V,N^{1/3}\}<v\le H.
\]
The restriction $4\theta<\nu$ ensures that the
approximation errors are smaller than the separation
between distinct rational products.

\begin{lemma}[Products of exceptional fractions]
\label{lem:products}
For each $b\in\mathcal A(w)$, choose an exceptional
fraction $s_b=u_b/v_b$ with
$\max\{V,N^{1/3}\}<v_b\le H$. Put
\[
 X=\{s_b:b\in\mathcal A(w)\},\qquad M=|X|.
\]
Then
\begin{equation}\label{eq:multiplicity}
 |\mathcal A(w)|\ll M,
\end{equation}
with an absolute implied constant. Moreover, the
product set $XX=\{ss':s,s'\in X\}$ satisfies
\begin{equation}\label{eq:products}
 |XX|\le4N.
\end{equation}
\end{lemma}

\begin{proof}
We first bound the number of indices assigned to a
single fraction $s=u/v\in X$. By
\eqref{eq:exception}, the values $\Phi(R(b))$
corresponding to these indices lie in an interval
of length $O(v^{-\nu})$ about $\Phi(s)$, since
$|\Phi'|\le(\log2)^{-1}$ on $[1,2]$.
The rotation estimate \eqref{eq:grid} therefore
gives at most
\[
 O(Nv^{-\nu}+1)=O(1)
\]
such indices, where we used $v>N^{1/3}$ and
$\nu>4$. This proves \eqref{eq:multiplicity}.

For $b,c\in[N,2N)$, write $\ell=b+c$. Then
\[
 R(b)R(c)=2^{a(b)+a(c)}3^{-\ell}.
\]
Since $0<a(d)-d\alpha<1$ for every $d\in\N$,
we have
\[
 \alpha\ell<a(b)+a(c)<\alpha\ell+2.
\]
Thus, for fixed $\ell$, the exponent $a(b)+a(c)$
has at most two possible values. As
$2N\le\ell\le4N-2$, the products $R(b)R(c)$
take at most $4N$ distinct values.

We transfer this bound to the products of elements
of $X$. For $b,c\in\mathcal A(w)$, all four
fractions lie in $[1,2]$, and hence
\[
 \begin{aligned}
 |s_bs_c-R(b)R(c)|
 &\le2|s_b-R(b)|+2|s_c-R(c)|\\
 &<\frac14V^{-\nu}.
 \end{aligned}
\]
Consequently, if
$R(b)R(c)=R(b')R(c')$, then
\[
 |s_bs_c-s_{b'}s_{c'}|<\frac12V^{-\nu}.
\]
Each rational product has reduced denominator at
most $H^2$. Distinct such products differ by at
least $H^{-4}$. Since
\[
 V^{-\nu}<V^{-4\theta}=H^{-4},
\]
the corresponding rational products must be equal.
Thus every value of $R(b)R(c)$ determines at most
one value of $s_bs_c$, proving \eqref{eq:products}.
\end{proof}

\subsection{Classes modulo rational powers}

We use \eqref{eq:products} to bound the number of
classes occupied by $X$ modulo rational powers.
A consequence of the Pl\"unnecke--Ruzsa inequality
controls the number of classes, while a lower bound
for linear forms in logarithms and a $3$-adic
valuation bound their sizes.

\begin{lemma}[Number of classes]\label{lem:classes}
Let $A$ be a finite nonempty subset of a torsion-free
abelian group $G$, written additively, and put
$L=|A+A|/|A|$. For every $q=2^k$, $k\ge1$,
the image of $A$ in $G/qG$ has at most $L^{3k}$
elements.
\end{lemma}

\begin{proof}
For finite subsets $A,B\subset G$, write
$A\pm B=\{a\pm b:a\in A,\ b\in B\}$ and
$m\cdot A=\{ma:a\in A\}$. Choose one representative
$a\in A$ from each class in the image of $A$ in
$G/qG$. The sets $a-q\cdot A$ belong to distinct
cosets of $qG$ and are therefore disjoint.
Since dilation by $q$ is injective, each contains
$|A|$ elements. Thus the number of classes is
at most
\[
 \frac{|A-q\cdot A|}{|A|}.
\]

The Pl\"unnecke--Ruzsa inequality
\cite[Theorem~1.2]{Petridis} gives
\[
 |A-2\cdot A|\le|A-A-A|\le L^3|A|.
\]
To iterate this bound, we use Ruzsa's triangle
inequality
\[
 |U-W|\,|V|\le|U-V|\,|V-W|
\]
for finite nonempty sets
\cite[proof of Theorem~1.2, (3)]{Petridis}.
Taking $U=A$, $V=2^{i-1}\cdot A$, and
$W=2^i\cdot A$, and using injectivity of
dilation, we obtain
\[
 \frac{|A-2^i\cdot A|}{|A|}
 \le
 \frac{|A-2^{i-1}\cdot A|}{|A|}
 \frac{|A-2\cdot A|}{|A|}
 \qquad(i\ge2).
\]
Induction yields
$|A-2^k\cdot A|\le L^{3k}|A|$, as required.
\end{proof}

For a positive rational number $r/s$ in lowest
terms, write
\[
 h(r/s)=\log\max\{r,s\}.
\]
Two positive rationals belong to the same class
modulo $q$th powers if their quotient is a $q$th
power in $\Q_{>0}$.

\begin{lemma}[Size of a class]\label{lem:one-class}
Let $4<\nu\le5$, $1<\theta<\nu/4$,
$N\in\DD$, and $w\ge1$. Let $X$ be as in
Lemma~\ref{lem:products}, with one representative
index chosen for each element of $X$.
Put $x=\log(N/w)$, and let $q$ be the least
power of two with $q\ge x^2$.

There exists an absolute constant $x_0>0$ such
that, whenever $x\ge x_0$, any two elements
$s_b,s_c\in X$ in the same class modulo rational
$q$th powers satisfy
\[
 \frac{s_c}{s_b}=\frac{R(c)}{R(b)},
 \qquad
 |c-b|\log3\le2w+\log2.
\]
Consequently, each class contains $O(w)$
elements of $X$, with an absolute implied constant.
\end{lemma}

\begin{proof}
It suffices to consider distinct elements
$s_b,s_c$ in the same class. Write
$s_c/s_b=g^q$, where $g\in\Q_{>0}\setminus\{1\}$.
Since each selected fraction has denominator at
most $H=e^w$ and numerator at most $2H$, the
reduced numerator and denominator of $s_c/s_b$
are at most $2H^2$. Hence
\begin{equation}\label{eq:ratioheight}
 qh(g)=h(s_c/s_b)\le2w+\log2\le3w.
\end{equation}
In particular, the height of $g$ is reduced by
a factor $q$.

We first prove that the rational-power relation
is an exact identity between the corresponding
prescribed fractions. Consider
\[
 \Lambda=(a(c)-a(b))\log2
 -(c-b)\log3-q\log g.
\]
Since $s_c/s_b=g^q$, we have
\[
 \Lambda=
 \log\frac{R(c)}{s_c}
 -\log\frac{R(b)}{s_b}.
\]
The logarithm has derivative at most one on
$[1,2]$. Applying \eqref{eq:exception} at
$b$ and $c$ therefore gives
\begin{equation}\label{eq:logupper}
 |\Lambda|<\frac18V^{-\nu}
 =\frac18\exp(-\nu w/\theta).
\end{equation}

Suppose that $\Lambda\ne0$. We apply Matveev's
lower bound
\cite[Corollary~2.3, (2.6), with (1.3)]{Matveev}.
For positive rational bases
$\gamma_i\ne1$ and integer coefficients $m_i$,
$1\le i\le3$, choose
\[
 A_i\ge\max\{h(\gamma_i),|\log\gamma_i|,0.16\},
 \qquad
 B=\max\left\{1,\max_{1\le i\le3}
 \frac{|m_i|A_i}{A_3}\right\}.
\]
For a nonzero linear form in their logarithms,
the cited bound gives
\[
 \log\left|\sum_{i=1}^3m_i\log\gamma_i\right|
 >-C_0A_1A_2A_3(1+\log B),
\]
where $C_0$ is absolute. The result does not
require multiplicative independence of the bases.

Apply this estimate with
\[
 (\gamma_1,\gamma_2,\gamma_3)=(2,3,g),
 \qquad
 (A_1,A_2,A_3)=(\log2,\log3,3w/q).
\]
These height parameters are admissible:
$|\log g|\le h(g)$, while
\eqref{eq:ratioheight} and $g\ne1$ give
\[
 \log2\le h(g)\le3w/q.
\]
Moreover, $b,c\in[N,2N)$ imply
\[
 |c-b|<N,\qquad
 |a(c)-a(b)|\le\alpha N+1.
\]
Thus the coefficients of $\log2$ and $\log3$
are $O(N)$, whereas that of $\log g$ has
absolute value $q$. Since $N/w=e^x$, the
weighted definition of $B$ yields
\[
 B\ll\frac{qN}{w},
 \qquad
 1+\log B\ll1+\log q+x.
\]
As $x^2\le q<2x^2$, Matveev's bound now gives
\[
 -\log|\Lambda|
 \ll\frac{w}{q}(1+\log q+x)
 \ll\frac wx,
\]
with absolute implied constants.

On the other hand, \eqref{eq:logupper} gives
\[
 -\log|\Lambda|>
 \frac{\nu w}{\theta}+\log8>4w.
\]
For $x$ sufficiently large, these estimates
are incompatible. Hence $\Lambda=0$, and
therefore
\[
 \frac{s_c}{s_b}
 =2^{a(c)-a(b)}3^{b-c}
 =\frac{R(c)}{R(b)}.
\]

Taking $3$-adic valuations, we obtain
$v_3(s_c/s_b)=b-c$. Since
$|v_3(z)|\log3\le h(z)$ for every positive
rational $z$, \eqref{eq:ratioheight} gives
\[
 |c-b|\log3\le h(s_c/s_b)
 \le2w+\log2.
\]
Fixing one representative index in a class,
all other representative indices therefore
lie in an interval of length $O(w)$.
They are distinct integers, so the class
contains $O(w)$ elements. The same bound
holds for a singleton, since $w\ge1$.
\end{proof}

\begin{lemma}[Count in a denominator range]
\label{lem:block}
There exist absolute constants $c,C,x_0>0$ such
that, for every $4<\nu\le5$, $1<\theta<\nu/4$,
$N\in\DD$, and $w\ge1$ with
$x=\log(N/w)\ge x_0$,
\[
 \frac{|\mathcal A(w)|}{N}
 \le C\exp\left(-\frac{cx}{\log(2+x)}\right).
\]
\end{lemma}

\begin{proof}
We may assume $\mathcal A(w)\ne\varnothing$.
Take $x_0\ge2$ large enough for
Lemma~\ref{lem:one-class}, and let $q=2^k$
be the least power of two with $q\ge x^2$.

By Lemma~\ref{lem:classes}, applied to $X$
in the multiplicative group $\Q_{>0}$, the
number of classes modulo rational $q$th
powers meeting $X$ is at most
\[
 \left(\frac{|XX|}{M}\right)^{3k}
 \le\left(\frac{4N}{M}\right)^{3k},
\]
where we used \eqref{eq:products}.
Lemma~\ref{lem:one-class} bounds the number
of elements in each class by $C_1w$, with
$C_1$ absolute. Consequently,
\[
 M\le C_1w\left(\frac{4N}{M}\right)^{3k}.
\]
Putting $y=M/(4N)$ and using $w/N=e^{-x}$,
we obtain
\[
 y^{3k+1}\le\frac{C_1}{4}e^{-x}.
\]
For sufficiently large $x$, the right-hand
side is at most $e^{-x/2}$. Moreover,
$q<2x^2$ gives
\[
 3k+1\ll\log(2+x),
\]
with an absolute implied constant. Taking
the $(3k+1)$st root therefore yields
\[
 \frac{M}{N}
 \ll\exp\left(-\frac{cx}{\log(2+x)}\right)
\]
for some absolute $c>0$. The conclusion
follows from \eqref{eq:multiplicity}, after
adjusting the absolute constants.
\end{proof}

\subsection{Summation over denominator ranges}

It remains to sum the preceding estimates.
The height bound \eqref{eq:height} ensures
that all exceptional denominator ranges
occur at logarithmic heights $o(N)$.
Consequently, the estimates of
Lemma~\ref{lem:block} apply in the tail
of a convergent series.


\begin{proof}[Proof of Proposition~\ref{prop:count}]
Fix $4<\nu\le5$ and choose
$1<\theta<\nu/4$. For $N\in\DD$, put
\[
 w_j=N\theta^{-j},\qquad
 x_j=\log(N/w_j)=j\log\theta
 \qquad(j=0,1,\ldots).
\]
The intervals $(w_j/\theta,w_j]$ partition
$(0,N]$. By \eqref{eq:height}, they cover
the logarithms of all exceptional
denominators $v>N^{1/3}$ once $N$ is
sufficiently large.

Suppose $\mathcal A(w_j)\ne\varnothing$.
Some exceptional denominator in that range
satisfies
\[
 \frac13\log N<\log v\le w_j
 <\theta\log v\le\theta W_N.
\]
In particular, $w_j\ge1$ for sufficiently
large $N$. Moreover, if $W_N>0$, then
\[
 x_j>\log\frac{N}{\theta W_N}
 \longrightarrow\infty
 \qquad(N\to\infty,\ N\in\DD),
\]
by \eqref{eq:height}. Thus every nonempty
range eventually satisfies the hypotheses
of Lemma~\ref{lem:block}. If $W_N=0$,
there are no exceptional fractions to count.

For fixed $\theta>1$, the series
\[
 \sum_{j\ge0}
 \exp\left(-\frac{cx_j}{\log(2+x_j)}\right)
\]
converges. Indeed, since $x_j=j\log\theta$,
its terms are bounded by
$\exp(-c_\theta\sqrt j)$ for all sufficiently
large $j$, with some $c_\theta>0$.

Given any $J\ge1$, every nonempty range
has $j\ge J$ for all sufficiently large $N$.
Lemma~\ref{lem:block} therefore gives
\[
 \limsup_{\substack{N\to\infty\\N\in\DD}}
 \frac1N\sum_{j\ge0}|\mathcal A(w_j)|
 \le C\sum_{j\ge J}
 \exp\left(-\frac{cx_j}{\log(2+x_j)}\right),
\]
where $C$ is absolute. Letting $J\to\infty$
shows that
\[
 \sum_{j\ge0}|\mathcal A(w_j)|=o_\nu(N).
\]
Every index admitting an exceptional
denominator greater than $N^{1/3}$
belongs to at least one of these sets.
Consequently,
\[
 E_N(N^{1/3},\nu)=o_\nu(N).
\]

Finally, Lemma~\ref{lem:small} gives
\[
 E_N(D,\nu)
 \le C\bigl(ND^{2-\nu}+N^{2/3}\bigr)
 +o_\nu(N),
\]
with an absolute constant $C$.
Dividing by $N$ and taking the upper
limit along $\DD$ proves the proposition.
\end{proof}


\section*{AI use declaration}

The author used ChatGPT and Codex in developing the counting
argument of Section~\ref{sec:count}, including the reduction to twofold
products and the use of rational-power classes, and in
formulating the diagonal construction in Section~\ref{sec:main}.
These tools were also used to review the exposition and
check proofs and references. The author takes full
responsibility for the mathematical arguments and citations.

\begin{thebibliography}{99}

\bibitem{Bugeaud}
Y. Bugeaud, \emph{Approximation by Algebraic Numbers}, Cambridge
Tracts in Mathematics \textbf{160}, Cambridge University Press,
Cambridge, 2004.

\bibitem{BN}
Y. Bugeaud and K. D. Nguyen,
\href{https://doi.org/10.2140/pjm.2023.326.17}{Some arithmetical
properties of convergents to algebraic numbers},
\emph{Pacific J. Math.} \textbf{326} (2023), no.~1, 17--36.

\bibitem{EM}
P. Erd\H{o}s and K. Mahler,
\href{https://doi.org/10.1112/jlms/s1-14.1.12}{Some arithmetical
properties of the convergents of a continued fraction},
\emph{J. London Math. Soc.} \textbf{14} (1939), no.~1, 12--18.

\bibitem{ErdosReminiscences}
P. Erd\H{o}s,
\href{https://old.renyi.hu/~p_erdos/1989-34.pdf}
{Some personal and mathematical reminiscences
of Kurt Mahler},
\emph{Austral. Math. Soc. Gaz.}
\textbf{16} (1989), no.~1, 1--2.

\bibitem{Fraenkel62}
A. S. Fraenkel,
\href{https://doi.org/10.1090/S0002-9947-1962-0144852-2}{On a
theorem of D. Ridout in the theory of Diophantine approximations},
\emph{Trans. Amer. Math. Soc.} \textbf{105} (1962), 84--101.

\bibitem{Fraenkel}
\bysame,
\href{https://doi.org/10.1112/jlms/s1-39.1.405}{Transcendental
numbers and a conjecture of Erd\H{o}s and Mahler},
\emph{J. London Math. Soc.} \textbf{39} (1964), 405--416.

\bibitem{FraenkelBorosh}
A. S. Fraenkel and I. Borosh,
\href{https://doi.org/10.1112/plms/s3-15.1.458}{Fractional dimension
of a set of transcendental numbers}, \emph{Proc. London Math. Soc.}
(3) \textbf{15} (1965), 458--470. Corrigendum, \emph{ibid.}
(3) \textbf{16} (1966), 192.

\bibitem{LelisMarques}
J. Lelis and D. Marques,
\href{https://doi.org/10.1017/S0004972716000770}
{On a problem of Erd\H{o}s and Mahler concerning
continued fractions},
\emph{Bull. Aust. Math. Soc.}
\textbf{95} (2017), no.~2, 183--186.

\bibitem{Matveev}
E. M. Matveev,
\href{https://doi.org/10.1070/IM2000v064n06ABEH000314}{An explicit
lower bound for a homogeneous rational linear form in the logarithms
of algebraic numbers. II}, \emph{Izv. Math.} \textbf{64} (2000),
no.~6, 1217--1269.

\bibitem{Petridis}
G. Petridis,
\href{https://doi.org/10.1007/s00493-012-2818-5}{New proofs of
Pl\"unnecke-type estimates for product sets in groups},
\emph{Combinatorica} \textbf{32} (2012), no.~6, 721--733,
\href{https://arxiv.org/abs/1101.3507v3}{arXiv:1101.3507v3}.

\bibitem{Ridout}
D. Ridout,
\href{https://doi.org/10.1112/S0025579300001182}{Rational
approximations to algebraic numbers},
\emph{Mathematika} \textbf{4} (1957), 125--131.

\bibitem{Schlickewei}
H. P. Schlickewei,
\href{https://www.numdam.org/item/CM_1992__82_3_245_0/}{The
quantitative Subspace Theorem for number fields},
\emph{Compositio Math.} \textbf{82} (1992), no.~3, 245--273.

\end{thebibliography}
\end{document}
